\documentclass[a4paper,12pt]{article}
\usepackage[utf8]{inputenc}
\usepackage[T1]{fontenc}
\usepackage[margin=2.5cm]{geometry}
\usepackage{ctex} % 支持中文
\usepackage{booktabs} % 优化表格线
\usepackage{array} % 表格列格式
\usepackage{graphicx} % 图片支持
\usepackage{titlesec} % 标题格式
\usepackage{xcolor} % 颜色支持
\usepackage{enumitem} % 列表优化

% 设置标题格式
\titleformat{\section}{\Large\bfseries\color{blue!40!black}}{}{0em}{}
\titleformat{\subsection}{\large\bfseries\color{gray!80!black}}{}{0em}{}

% 元数据
\title{\textbf{2024年度环境、社会及管治（ESG）报告节选}}
\author{公司董事会办公室}
\date{2024年3月}

\begin{document}

\maketitle

% 已移除摘要和目录，直接开始正文

\section{1. 前言：服务国家战略，守护自贸港门户}

2024年是自贸港封关运作的攻坚之年。作为深耕本省的航空基础设施与临空产业开发平台，本公司肩负着“外防输入、内保畅通”的重要使命。我们紧扣“一主一次，两支一货运”的机场群布局，通过提升枢纽能级，全力服务于国家开放战略。

在这一年里，我们不仅关注飞机的起降与旅客的输送，更将目光投向了更长远的未来——如何让蓝天更蓝，如何让旅客的笑容更灿烂，以及如何让企业的根基更稳固。

\section{2. 环境范畴（Environmental）：守护碧海蓝天}

我们深知，绿色是自贸港最鲜明的底色。2024年，公司积极响应国家“双碳”目标，通过技术创新与精细化管理，显著降低了运营过程中的碳足迹。

\subsection{2.1 碳排放与能源资源消耗实绩}

为了让公众清晰了解我们的环境影响，我们将2024年度的核心环境数据整理如下。我们坚持“每一度电都可追溯，每一滴水都物尽其用”的原则。

\begin{table}[h]
    \centering
    \caption{2024年度温室气体排放与能源消耗数据表}
    \label{tab:env_data}
    \renewcommand{\arraystretch}{1.3}
    \begin{tabular}{p{6cm} p{4cm} p{4cm}}
        \toprule
        \textbf{指标名称} & \textbf{数值} & \textbf{单位} \\
        \midrule
        \textbf{温室气体排放总量} & \textbf{51,518.71} & \textbf{吨CO$_2$当量} \\
        \quad - 范围1（直接排放，如汽油燃烧） & 11,214.39 & 吨CO$_2$当量 \\
        \quad - 范围2（间接排放，如外购电力） & 40,304.32 & 吨CO$_2$当量 \\
        碳排放强度 & 0.12 & 吨CO$_2$当量/万元 \\
        \midrule
        \textbf{能源消耗} & & \\
        综合能源消耗量 & 13,538.22 & 吨标准煤 \\
        外购电力 & 7,511.05 & 万千瓦时 \\
        汽油消耗 & 19.10 & 万升 \\
        \midrule
        \textbf{水资源与废弃物} & & \\
        水资源消耗总量 & 145.94 & 万吨 \\
        中水回用量（循环利用水） & 13.61 & 万吨 \\
        无害废弃物处理量 & 6,230.58 & 吨 \\
        合规处理率 & 100 & \% \\
        \bottomrule
    \end{tabular}
\end{table}

\subsection{2.2 打造“双碳机场”：让飞机“插电”休息}

传统的机场运行中，飞机停在地面时如果通过燃烧航空燃油来维持空调和照明（即使用APU辅助动力装置），会产生巨大的噪音和碳排放。为了解决这个问题，我们在旗下核心门户机场大力推广了“APU替代设施”。

简单来说，这就像是让飞机在地面“插电”休息。通过使用地面的电力设备直接为飞机供电供气，飞机不再需要消耗燃油。这项技术的应用效果显著：

\begin{itemize}
    \item \textbf{设备投入}：已投入使用31套APU替代设备（包括近机位18套，远机位13套）。
    \item \textbf{减碳成效}：全年累计减少碳排放 \textbf{25,382} 吨CO$_2$当量。
    \item \textbf{行业认可}：旗下核心门户机场因此荣获“双碳机场”二星级称号。
\end{itemize}

\subsection{2.3 清洁能源与绿色建筑：从屋顶到地面的绿色革命}

除了减少排放，我们还努力“生产”绿色能源，并将绿色理念融入建筑设计中。

\begin{itemize}
    \item \textbf{光伏发电}：
    \begin{itemize}
        \item \textbf{商业体项目}：旗下核心商业综合体利用屋顶建设光伏电站，年均发电量达 \textbf{450万} 千瓦时，这满足了该商业体约10\%的总用电需求，每年减少碳排放 2,263.59 吨。
        \item \textbf{临空产业园}：屋顶光伏装机容量达到 8.6MWp，预计每年可减少碳排放 25.75 万吨。
    \end{itemize}
    
    \item \textbf{新能源车辆应用}：
    \begin{itemize}
        \item 在核心机场，我们大力推进“油改电”。目前新能源车占比已达 \textbf{43.6\%}（共计221辆）。
        \item 配套新建充电桩约 100 个，全年通过车辆电动化减少碳排放 761 吨。
    \end{itemize}

    \item \textbf{绿色建筑认证}：
    \begin{itemize}
        \item 在建的省会城市地标项目“海南中心”（双子塔）已获得 \textbf{LEED CS 金级预认证}。
        \item 某新区F07项目获得绿色建筑一星级标准认证。
    \end{itemize}
\end{itemize}

\section{3. 社会范畴（Social）：共建美好空港}

机场不仅仅是交通枢纽，更是连接人与人的温情纽带。2024年，我们在保障安全、提升服务、关爱员工和回馈社会方面交出了务实的答卷。

\subsection{3.1 运营绩效与服务创新}

随着旅游市场的强劲复苏，我们的业务量实现了稳步增长。

\begin{table}[h]
    \centering
    \caption{2024年度主要运营数据概览}
    \label{tab:ops_data}
    \renewcommand{\arraystretch}{1.3}
    \begin{tabular}{p{8cm} p{4cm}}
        \toprule
        \textbf{运营指标} & \textbf{数据实绩} \\
        \midrule
        旗下控股及管理输出机场数量 & 9 家 \\
        总体起降架次 & 16.11 万架次 \\
        旅客吞吐量 & 2,462.62 万人次 \\
        核心门户机场排名 & 重回“2000万级”俱乐部（全国第29位） \\
        重点区域机场新航线 & 开通首条国际客运航线（至吉隆坡） \\
        \bottomrule
    \end{tabular}
\end{table}

为了让旅客的出行更加顺畅，我们推出了两项人性化的服务创新：

\begin{itemize}
    \item \textbf{通关“零等待”}：实施“一机双屏”海关安检融合查验模式。旅客只需接受一次检查，行李即可同时完成海关和安检的查验，大幅缩短了排队时间。
    \item \textbf{支付“无障碍”}：针对外籍旅客入境支付难的痛点，我们在核心机场启用了“一站式综合服务中心”。外籍友人在抵达的第一时间，就能搞定移动支付绑定、购买电话卡等事宜，真正实现了“落地即融入”。
\end{itemize}

\subsection{3.2 安全应急：直面超强台风“摩羯”的考验}

2024年9月，超强台风“摩羯”正面袭击本省。这不仅是对基础设施的考验，更是对我们应急体系的实战检验。

\begin{itemize}
    \item \textbf{严阵以待}：在台风登陆前，我们迅速启动了“防汛防风专项应急预案”。运管委提前进行航班动态调整，对机坪设施进行了地毯式的加固，确保“无死角”。
    \item \textbf{快速复航}：台风过境后，展现了惊人的“海南速度”。我们迅速清理跑道异物，修复受损设施。9月3日至7日期间，核心机场迅速恢复运行，累计执行航班 \textbf{1,300} 架次，有力保障了进出岛通道的畅通。
    \item \textbf{日常安全管理}：
    \begin{itemize}
        \item \textbf{数字化转型}：上线了“安全智控”平台。我们将厚厚的纸质手册变成了数字化的指令，平台关联制度手册 9,626 项，全年生成电子检查单 31 万份，确保安全标准落实到每一个动作。
        \item \textbf{安全成绩单}：全年一般征候万架次率为 0，责任事故为 0，安全培训覆盖率达到 100\%。
    \end{itemize}
\end{itemize}

\subsection{3.3 员工发展与多元包容}

员工是公司最宝贵的财富。我们致力于构建一个多元、平等且充满活力的工作环境。

\begin{table}[h]
    \centering
    \caption{2024年度人力资源数据表}
    \label{tab:hr_data}
    \renewcommand{\arraystretch}{1.3}
    \begin{tabular}{p{6cm} p{4cm}}
        \toprule
        \textbf{指标项目} & \textbf{数据统计} \\
        \midrule
        员工总数 & 10,611 人 \\
        女性员工占比 & 38.9\% (4,133人) \\
        少数民族员工 & 179 人 \\
        培训总投入 & 1,119.54 万元 \\
        人均培训时长 & 105.88 小时 \\
        董事会女性占比 & 22\% (2/9) \\
        \bottomrule
    \end{tabular}
\end{table}

\subsection{3.4 社会贡献：和合共生}

我们不仅追求商业成功，更不忘回馈社会。

\begin{itemize}
    \item \textbf{乡村振兴}：全年投入乡村振兴资金 \textbf{190 万元}，重点改善对口帮扶村的基础设施与教育条件。
    \item \textbf{志愿服务}：我们的“红马甲”活跃在各个角落。全年参与志愿服务超过 10,000 人次，总时长累计达到 \textbf{500,000 小时}。
    \item \textbf{重大保障}：圆满完成了博鳌亚洲论坛年会、消博会、少数民族运动会等国家级重大活动的保障任务，展示了“中国服务”的良好形象。
\end{itemize}

\section{4. 管治范畴（Governance）：筑牢合规基石}

透明、高效的公司治理是企业可持续发展的基石。我们建立了完善的治理架构，确保权力在阳光下运行。

\subsection{4.1 治理结构与投资者沟通}

我们严格遵守上市公司监管要求，构建了规范的“五会一层”治理体系。

\begin{itemize}
    \item \textbf{董事会构成}：董事会成员共9人，其中独立董事3人，确保了决策的客观性与独立性。
    \item \textbf{会议与披露}：全年共召开股东大会5次，董事会13次。发布各类公告70份，确保所有投资者都能及时、公平地获取公司信息。
\end{itemize}

\subsection{4.2 合规管理与反腐败}

为了防范风险，我们构建了“三全三化”大监督体系，将风险控制在萌芽状态。

\begin{table}[h]
    \centering
    \caption{2024年度合规与风控数据表}
    \label{tab:gov_data}
    \renewcommand{\arraystretch}{1.3}
    \begin{tabular}{p{8cm} p{4cm}}
        \toprule
        \textbf{工作项目} & \textbf{执行情况} \\
        \midrule
        梳理风险点 & 358 个 \\
        反腐败培训覆盖率（管理层） & 95\% \\
        反腐败培训覆盖率（员工） & 90\% \\
        反腐败培训总时长 & 400 小时 \\
        供应商廉洁承诺书签署率 & 100\% \\
        \bottomrule
    \end{tabular}
\end{table}

\subsection{4.3 经济绩效延伸}

作为自贸港建设的重要参与者，我们积极布局离岛免税业务，促进消费回流。2024年，公司间接参与的5家离岛免税店线下销售额约为 \textbf{52 亿元}，市场占比约 \textbf{17\%}，为地方经济发展做出了实质性贡献。

\end{document}